% ABSTRACT — exactly one page
\newpage
\thispagestyle{plain}
\begingroup
\setlength{\parskip}{6pt}
\setstretch{1.5}
\fontsize{12}{18}\selectfont

\begin{center}
    {\fontsize{16}{24}\selectfont\bfseries ABSTRACT\\[12pt]}
\end{center}

Large Language Models (LLMs) can translate plain-English questions into SQL, but a plain
Natural Language to SQL (NL2SQL) pipeline breaks down inside an organization. Business
rules such as ``revenue excludes cancelled orders'' are not in the schema, every session
starts from zero, large schemas do not fit in a prompt, and executing model-written SQL on a
live database is risky.

This project presents the \textbf{Context Memory Framework (CMF)}, a context-aware NL2SQL
assistant built around a \textbf{persistent, governed context-memory layer}. CMF stores
semantic, procedural and episodic memories plus structured working memory, keeps extracted
memories \textit{pending} until a person approves them, versions and audits every change,
and never stores one-time overrides as permanent rules. Memories, schema tables and past
questions are indexed in \textbf{Qdrant} and retrieved by hybrid dense--sparse search with
reciprocal rank fusion, a weighted first-stage score and a cross-encoder reranker, always
after authorization filters. A token-budgeted context builder assembles a prioritized prompt
that frames memories as untrusted data, and a deterministic \texttt{sqlglot} guard allows
only one read query, or one single-table change that an admin must approve.

The system is implemented with FastAPI, SQLAlchemy, Google Gemini and a server-rendered
Jinja2/htmx interface. All 66 automated tests pass. On a 120-case synthetic benchmark, the
current pipeline reaches Recall@3 of 0.79 and MRR of 0.73, includes the target business
context in 97.5\% of built prompts, and shows no cross-tenant or cross-user memory leakage.

\vspace{0.6cm}
\noindent\textbf{Keywords:}\enspace Natural Language to SQL, Large Language Models,
Context Engineering, Long-Term Memory, Retrieval-Augmented Generation, Hybrid Search,
Reranking, SQL Guardrails, FastAPI.

\endgroup
\clearpage
